\textbf{Phần 1. Trắc nghiệm lựa chọn}
\setcounterpageref{ex}{0}

\begin{ex} %[2D1N01]
Số gồm $5$ chục và $7$ đơn vị viết là
\choice
{$75$}
{$507$}
{$50$}
{\True $57$}
\loigiai{
Số gồm $5$ chục và $7$ đơn vị là $57$.\\
\textbf{Đáp án của câu hỏi là D}
}
\end{ex}

\begin{ex} %[2D1N01]
Số lớn nhất có hai chữ số là
\choice
{\True $99$}
{$90$}
{$100$}
{$98$}
\loigiai{
Số lớn nhất có hai chữ số là $99$.\\
\textbf{Đáp án của câu hỏi là A}
}
\end{ex}

\begin{ex} %[2D1H01]
Số bé nhất có hai chữ số giống nhau là
\choice
{$10$}
{\True $11$}
{$22$}
{$00$}
\loigiai{
Số bé nhất có hai chữ số giống nhau là $11$.\\
\textbf{Đáp án của câu hỏi là B}
}
\end{ex}

\begin{ex} %[2D1H01]
Số $84$ đọc là
\choice
{\True Tám mươi tư}
{Tám mươi bốn}
{Tám bốn}
{Bốn mươi tám}
\loigiai{
Số $84$ đọc là Tám mươi tư (hoặc tám mươi bốn).\\
\textbf{Đáp án của câu hỏi là A}
}
\end{ex}

\begin{ex} %[2D1H01]
Trong số $68$, chữ số 6 có giá trị là
\choice
{$6$}
{$68$}
{\True $60$}
{$8$}
\loigiai{
Chữ số 6 ở hàng chục nên có giá trị là $60$.\\
\textbf{Đáp án của câu hỏi là C}
}
\end{ex}

\begin{ex} %[2D1H01]
Điền dấu thích hợp: $45 \dots 54$.
\choice
{$>$}
{$=$}
{Không so sánh được}
{\True $<$}
\loigiai{
Hàng chục: $4 < 5$ nên $45 < 54$.\\
\textbf{Đáp án của câu hỏi là D}
}
\end{ex}

\begin{ex} %[2D1V01]
Sắp xếp các số $35, 53, 30, 50$ theo thứ tự từ bé đến lớn:
\choice
{$53, 50, 35, 30$}
{\True $30, 35, 50, 53$}
{$30, 50, 35, 53$}
{$35, 30, 50, 53$}
\loigiai{
Thứ tự từ bé đến lớn: $30 < 35 < 50 < 53$.\\
\textbf{Đáp án của câu hỏi là B}
}
\end{ex}

\begin{ex} %[2D1V01]
Từ 2 chữ số $3$ và $8$ lập được tất cả bao nhiêu số có hai chữ số khác nhau?
\choice
{$1$ số}
{$3$ số}
{\True $2$ số ($38$ và $83$)}
{$4$ số}
\loigiai{
Các số là $38$ và $83$ (có 2 số).\\
\textbf{Đáp án của câu hỏi là C}
}
\end{ex}

\begin{ex} %[2D1V01]
Viết số $76$ thành tổng của chục và đơn vị:
\choice
{\True $70 + 6$}
{$7 + 6$}
{$70 + 60$}
{$76 + 0$}
\loigiai{
$76 = 70 + 6$.\\
\textbf{Đáp án của câu hỏi là A}
}
\end{ex}

\begin{ex} %[2D1V01]
Số tròn chục liền sau số $40$ là
\choice
{$41$}
{\True $50$}
{$30$}
{$49$}
\loigiai{
Số tròn chục liền sau của $40$ là $50$.\\
\textbf{Đáp án của câu hỏi là B}
}
\end{ex}

\begin{ex} %[2D1V01]
Có bao nhiêu số tròn chục có hai chữ số?
\choice
{$8$ số}
{$10$ số}
{\True $9$ số (từ $10$ đến $90$)}
{$11$ số}
\loigiai{
Các số là $10, 20, 30, 40, 50, 60, 70, 80, 90$ (có 9 số).\\
\textbf{Đáp án của câu hỏi là C}
}
\end{ex}

\begin{ex} %[2D1C01]
Tìm một số có hai chữ số biết chữ số hàng chục là số lớn nhất có 1 chữ số, chữ số hàng đơn vị kém chữ số hàng chục 5 đơn vị.
\choice
{$95$}
{$83$}
{$93$}
{\True $94$}
\loigiai{
Hàng chục: 9. Hàng đơn vị: $9 - 5 = 4$. Số đó là $94$.\\
\textbf{Đáp án của câu hỏi là D}
}
\end{ex}

\textbf{Phần 2. Trắc nghiệm đúng sai}
\setcounterpageref{ex}{0}

\begin{ex} %[2D1H01]
Xét tính đúng sai của các nhận định về số có hai chữ số:
\choiceTF
{\True $10$ là số bé nhất có hai chữ số.}
{\True $99$ là số lớn nhất có hai chữ số.}
{Số $100$ là số có hai chữ số.}
{\True Số $100$ là số có ba chữ số.}
\loigiai{
\begin{itemize}
    \item $10$ là số bé nhất có hai chữ số: Đúng.
    \item $99$ là số lớn nhất có hai chữ số: Đúng.
    \item Số $100$ có 3 chữ số. Ý 3 Sai, Ý 4 Đúng.
\end{itemize}
\begin{itemize}
    \item a) \textbf{Đúng}.
    \item b) \textbf{Đúng}.
    \item c) \textbf{Sai}.
    \item d) \textbf{Đúng}.
\end{itemize}
}
\end{ex}

\begin{ex} %[2D1V01]
Cho các số: $48, 84, 40, 80$. Xét tính đúng sai:
\choiceTF
{\True Số lớn nhất trong các số trên là $84$.}
{\True Số bé nhất trong các số trên là $40$.}
{Có 3 số tròn chục trong 4 số trên.}
{\True Chỉ có 2 số tròn chục là $40$ và $80$.}
\loigiai{
\begin{itemize}
    \item Số lớn nhất là $84$: Đúng.
    \item Số bé nhất là $40$: Đúng.
    \item Có 2 số tròn chục ($40, 80$). Ý 3 Sai, Ý 4 Đúng.
\end{itemize}
\begin{itemize}
    \item a) \textbf{Đúng}.
    \item b) \textbf{Đúng}.
    \item c) \textbf{Sai}.
    \item d) \textbf{Đúng}.
\end{itemize}
}
\end{ex}

\textbf{Phần 3. Trả lời ngắn}
\setcounterpageref{ex}{0}

\begin{ex} %[2D1H01]
Viết số gồm $8$ chục và $3$ đơn vị.
\shortans{83}
\loigiai{
\textbf{Đáp số:} $83$\\
83.
}
\end{ex}

\begin{ex} %[2D1H01]
Tìm số bé nhất có hai chữ số.
\shortans{10}
\loigiai{
\textbf{Đáp số:} $10$\\
10.
}
\end{ex}

\begin{ex} %[2D1V01]
Tìm số lớn nhất có hai chữ số khác nhau.
\shortans{98}
\loigiai{
\textbf{Đáp số:} $98$\\
98.
}
\end{ex}

\begin{ex} %[2D1V01]
Tìm số bé nhất có hai chữ số khác nhau.
\shortans{10}
\loigiai{
\textbf{Đáp số:} $10$\\
10.
}
\end{ex}

\begin{ex} %[2D1V01]
Tính tổng chữ số hàng chục và hàng đơn vị của số $45$.
\shortans{9}
\loigiai{
\textbf{Đáp số:} $9$\\
$4 + 5 = 9$.
}
\end{ex}

\begin{ex} %[2D1C01]
Tìm số tròn chục lớn nhất có hai chữ số mà tổng các chữ số của nó bằng 8.
\shortans{80}
\loigiai{
\textbf{Đáp số:} $80$\\
Số tròn chục có dạng $\overline{a0} \Rightarrow a + 0 = 8 \Rightarrow a = 8$. Số đó là $80$.
}
\end{ex}
